\section{Introduction}\label{sec:intro}

Rice is the backbone of agriculture in Northeast Thailand, where most of the crop is grown as a single, rainfed wet-season (main-season) crop on the sandy and partly saline soils of the Khorat Plateau. Yields in this environment depend on the timing and amount of monsoon rainfall, on floods along the Mun, Chi and Mekong river systems, on heat and humidity around flowering, and on how farmers allocate land, labour and fertiliser between rice and upland cash crops such as cassava, sugarcane and maize. Timely district-level yield forecasts are therefore valuable for food-security monitoring, for targeting drought and flood relief, for the design and settlement of crop insurance, and for market and price-support policy.

The data available for such forecasts have never been richer. Optical satellite time series describe canopy development \citep{tucker1979red,huete2002overview}, water-sensitive indices track flooding and transplanting in paddies \citep{gao1996ndwi,xiao2005mapping}, thermal sensors measure land-surface temperature \citep{wan2008new}, gridded rainfall and land reanalyses provide weather and soil moisture \citep{funk2015climate,munozsabater2021era5}, and cloud platforms make these archives accessible at the scale of administrative units \citep{gorelick2017google}. Machine-learning (ML) models that combine such predictors now routinely report high accuracy \citep{jeong2016random,cai2019integrating,khaki2019crop,vanklompenburg2020crop}.

Three methodological problems nevertheless limit what can be learned from this literature. The first is \emph{evaluation leakage}. Yield data form a spatio-temporal panel in which errors are correlated within seasons and among neighbouring units. Randomly partitioned cross-validation lets a model learn from future seasons and from neighbours in the same season, and seemingly innocuous predictors---contemporaneous prices, anomalies computed against full-period means, spatial averages of neighbours' outcomes---import information that would not be available at forecast time \citep{kaufman2012leakage,roberts2017cross,kapoor2023leakage}. The resulting optimism is well documented for time series and spatial ecology \citep{tashman2000out,bergmeir2012use,roberts2017cross} but is rarely quantified for multi-source yield models. The second is \emph{explanation without ground truth}. Feature attributions such as SHAP values \citep{lundberg2017unified,lundberg2020local} are increasingly used to rank the drivers of yield, yet with dependent predictors credit can be distributed in counter-intuitive ways \citep{aas2021explaining,hooker2021unrestricted}, and with real data the true drivers are unknown, so the faithfulness of such explanations cannot be checked. The third is \emph{incomplete integration}. Economic predictors---own and competing-crop prices that govern land and input allocation \citep{nerlove1956estimates}---and explicitly spatial predictors are seldom evaluated alongside weather and satellite data, and point forecasts are rarely accompanied by calibrated uncertainty \citep{vovk2005algorithmic,lei2018distribution}.

This study addresses these problems with a simulation benchmark. Following the logic of simulation studies for statistical methods \citep{morris2019using}, we construct a process-informed data-generating process (DGP) for main-season rice in \nDistricts{} districts of Northeast Thailand over \nSeasons{} seasons. The DGP encodes agronomic mechanisms (water stress at establishment and late in the season, floods, heat, humid-season disease, salinity), economic mechanisms (input intensity, competition from cassava, sugarcane and maize, and a harvest-to-price feedback), persistent management and local pest outbreaks. It emits observations that mimic real sources---monthly weather and soil moisture, cloud-contaminated satellite indices, national and local prices, and imperfect map layers---and it stores every additive yield component, so that model explanations can be compared with the truth. The same code ingests real data with an identical schema, so the benchmark doubles as a reproducible template for empirical work. We stress at the outset that all results reported here concern the synthetic benchmark; they are not empirical statements about Thai rice yields.

We ask four research questions and state hypotheses for each, formulated at the design stage:
\begin{description}
\item[RQ1 (accuracy).] How much do multi-source learners improve on persistence, climatology and an econometric panel model under leakage-safe evaluation? \textbf{H1}: multi-source learners are significantly more accurate than all baselines.
\item[RQ2 (evaluation).] How much do random cross-validation and subtle feature leakage overstate skill? \textbf{H2}: both produce materially optimistic error estimates.
\item[RQ3 (attribution).] Which feature groups carry predictive information, and do SHAP attributions recover the true mechanisms? \textbf{H3}: satellite features carry unique information (pest outbreaks invisible to weather data), whereas the moisture, water and satellite groups partly substitute for each other. \textbf{H4}: SHAP recovers strong, well-observed mechanisms but not weak or year-confounded ones such as price effects.
\item[RQ4 (timeliness and uncertainty).] How does skill evolve during the season, and are prediction intervals calibrated? \textbf{H5}: skill increases with the data cut-off month and the value of satellite data grows after panicle initiation. \textbf{H6}: conformal intervals are near-nominal on average but under-cover in shock seasons.
\end{description}

The paper makes five contributions: (i) an open, documented benchmark DGP with ground-truth yield components for multi-source yield forecasting; (ii) a quantification of how much common evaluation shortcuts inflate apparent skill; (iii) a ground-truth test of grouped SHAP attributions in the presence of dependent predictor groups; (iv) evidence on the in-season value of satellite data and on the calibration of distribution-free prediction intervals; and (v) a leakage-tested, fully reproducible code base with a documented pathway to real Thai data. Section~\ref{sec:related} reviews related work, Section~\ref{sec:design} describes the benchmark, Section~\ref{sec:methods} the methods and Section~\ref{sec:results} the results; Section~\ref{sec:discussion} discusses implications and limitations and Section~\ref{sec:conclusion} concludes.
